\section{Limits of this audit}
\label{sec:limits}

\begin{enumerate}
\item \textbf{The mathematics was not read.} This note establishes that a particular objection fails and that standard mechanical checks pass. It does not explain the proof and is not a substitute for mathematical review.
\item \textbf{The trust base remains.} The conclusion rests on the consistency of Lean's type theory with the three standard axioms, on the correctness of at least one of the checkers that accepted the export
(we ran one, \texttt{con-leche}, ourselves; others are reported on the verification page), on the faithfulness of \texttt{lean4export}, and on the adequacy of Mathlib's definitions of the Riemann zeta function and of \lean{RiemannHypothesis}.
None of these is questioned in \N, and we did not audit them beyond what Section~\ref{sec:defects} describes.
\item \textbf{The source-to-export build was not reproduced.} See Section~\ref{sec:defects}, exit~(iv).
\item \textbf{A dependency graph is syntactic.} It tells exactly which declarations a term refers to. Absence from the export means ``not used''; presence means ``referred to'', which is weaker than ``essential''.
The split between library and release is by name against the Mathlib environment and may misplace a small number of automatically generated auxiliary declarations.
\item \textbf{Counts measure extent, not depth.} That the bundle witness depends on \NpackageThm\ theorems shows that it is not obtained from nothing. It does not show that those theorems are deep; we did not grade them.
\item \textbf{Not every agent finding was re-verified.} Those that were not are marked \textsf{G} without \textsf{S} or \textsf{K} in Table~\ref{tab:facts}.
\item \textbf{The version of \N.} We reviewed the 17-page version with the hash given in Table~\ref{tab:prov}. \N\ has been revised more than once, and a later version may differ.
\end{enumerate}
